\chapter{Lernziele}

\section{Lernziele Kapitel 2: Einführung in Python und erste Schleifen}
\begin{todolist}
    \item Ich kann mit der Turtle einfache Formen zeichnen (Dreiecke, Rechtecke, Blitz, etc.)
    \item Ich kann mit einem breiten Stift ausgefüllte Rechtecke zeichnen (Türe beim Haus)
    \item Ich setze die \lstinline|for _ in range(zahl)|-Schleife ein, um sich wiederholende Tätigkeiten auszuführen.
    \item Ich kann mit der Turtle regelmässige Vielecke (Vierecke, Fünfecke, Sechsecke etc.) zeichnen.
    \item Ich kann ``Kreise'' als Vielecke mit grosser Eckenzahl zeichnen.
    \item Ich kann verschachtelte Schleifen verwenden (\lstinline|for _ in range(zahl)|-Schleifen innerhalb von \lstinline|for _ in range(zahl)|-Schleifen).
    %\item Ich kann Blütenblätter zeichnen, die aus Drittel, Viertel-, Fünftel-, usw.-Kreise bestehen.
    \item Ich kann Text in der Konsole mit \lstinline|print("...")| ausgeben.
    \item Ich setze die Operationen \lstinline|+| und \lstinline|*|, um komplizierte Texte (Zeichenketten) zu erzeugen.
    \item Ich kenne die Operationen auf Zahlen, die auf Seite 19 oben aufgelistet sind, und kann diese auch einsetzen.
    \item Ich kann innerhalb eines Programms eine Streckenlänge mit Pythagoras berechnen.
    \item Ich kann die Farbe und Breite des Turtle-Stifts verändern (\lstinline|t.color("...")|, \lstinline|t.width(...)|)
    \item Ich kann die Turtle bewegen, ohne zu zeichnen, indem ich den \enquote{Stift} ab- sowie aufsetze (\lstinline|t.pu()|, \lstinline|t.pd()|)
\end{todolist}
\section{Lernziele Kapitel 3: Variablen, Datentypen \& Debugging}
\begin{todolist}
    \item Ich kann neue Variablen erstellen, um die Resultate einfacher Berechnungen zu speichern
    \item Ich kann mit dem Gleich-Operator (\lstinline|=|) eine neue Variable erstellen und kann diese in einem Code sinnvoll verwenden.
    \item Ich kann Speicherinhalte ändern, indem ich die Befehle \lstinline|+=|, \lstinline|-=|, \lstinline|*=| und \lstinline|/=| nutze.
    \item Ich kann mithilfe des \lstinline|input|-Befehls neue Variablen während der Ausführung des Codes erstellen.
    \item Ich kenne den Unterschied zwischen der Division (\lstinline|/|) und der Ganzzahl-Division (\lstinline|//|) und kann beide Operatoren in geeigneten Situationen anwenden.
    \item Ich kann den Modulo-Operator \lstinline|%| verwenden, um den Rest einer Ganzzahl-Division zu berechnen.
    \item Ich kann Variablen nutzen, um eine Spirale oder ähnliche, sich stetig vergrössernde oder verkleinernde Formen zu zeichnen.
    \item Ich kann Zahlenfolgen (z.B. alle Quadratzahlen bis zu einem bestimmten Wert) mit Python erstellen.
\end{todolist}

\section{Lernziele Kapitel 4: Funktionen}
\begin{todolist}
    \item Ich kann eigene Befehle in Python definieren und einsetzen.
    \item Ich verstehe den Unterschied zwischen Haupt- und Unterprogramm.
    \item Ich kann Befehle mit Parameter definieren und einsetzen, wie z.B. \lstinline|quadrat(seite)|.
    \item Ich kann die Lebensdauer von Variablen (inkl. Parametern) beschreiben (lokal oder global).
    \item Ich kann beschreiben, wie sich ein Parameter von andern Variablen unterscheidet.
    \item Ich kann Befehle mit mehreren Parameter definieren und einsetzen, wie z.B. \lstinline|vieleck(ecken, laenge, farbe)|.
    \item Ich kann einen Befehl innerhalb eines anderen Befehls aufrufen und dabei Variablenwerte übergeben.
    \item Ich kann Speicherinhalte ändern, indem ich den Inhalt durch den Wert eines Ausdrucks überschreibe, z.B. \lstinline|seite = seite * 2 + 1|
    \item Ich kann einfache Befehle kombinieren, um komplexe Probleme zu lösen (Beispiel Häuserreihe).
    \item Ich kann Funktionen erstellen, die einen oder mehrere Parameter sowie einen \lstinline|return|-Wert haben.
    \item Ich verstehe den Unterschied zwischen einem \lstinline|print|-Befehl und einem \lstinline|return|-Befehl
    \item Ich kann das Resultat einer Funktion mit \lstinline|return| in einer Variable im Hauptprogramm speichern und diese weiterverwenden
    \item Ich kann Funktionen mit einem \lstinline|return|-Ausdruck erstellen, den Rückgabewert der Funktion im Hauptprogramm als Variable speichern und den Wert der Variable mit \lstinline|print| anzeigen lassen.
    \item Ich kann eine (Unter-)Funktion innerhalb einer anderen (Haupt-)Funktion aufrufen, und den Rückgabewert der Unterfunktion in der Hauptfunktion weiterverwenden.
    \item Ich nutze den modularen Programmentwurf, um komplexe Probleme in Teilprobleme aufzuteilen.
    \item Ich kann den \lstinline|return|-Wert einer Funktion direkt und ohne Zwischenspeicherung in einem Ausdruck verwenden, z.B. \lstinline|if rechteck_flaeche(19, 12) < 100: ...|.
\end{todolist}

\section{Lernziele Kapitel 5: Verzweigungen und Logische Ausdrücke}
\begin{todolist}
    \item Ich kann die \lstinline|if|-Anweisung einsetzen und verwende dabei die Vergleichsoperatoren \lstinline|==|, \lstinline|!=|, \lstinline|>|, \lstinline|<|, \lstinline|<=| und \lstinline|>=|.
    \item Ich kann die \lstinline|if|-\lstinline|else|-Struktur für Verzweigungen mit zwei Fällen verwenden.
    \item Ich kann die \lstinline|if|-\lstinline|elif|-\lstinline|else|-Anweisung für Verzweigungen mit beliebig vielen Fällen verwenden.
    \item Ich kenne den Unterschied im Programmablauf zwischen der \lstinline|if|-\lstinline|elif|-\lstinline|else|- und der \lstinline|if|-\lstinline|if|-\lstinline|else|-Struktur.
    \item Ich kann die logischen Operatoren \lstinline|and|, \lstinline|or| sowie \lstinline|not| einsetzen, um boolsche Ausdrucke zu kombinieren.
    \item Ich kann die \lstinline|break|-Anweisung einsetzen, um Schleifen unter bestimmten Bedingungen abzubrechen.
    \item Ich kann \lstinline|while|-Schleifen einsetzen, um eine Schleife solange auszuführen wie eine Bedingung wahr ist.
    \item Ich kann erklären, weshalb ein Code der nach einem ausgeführten \lstinline|for _ in range(zahl)|-Ausdruck steht, nicht mehr ausgeführt wird.
    \item Ich kann für ein einfaches, gegebenes Code-Beispiel bestimmen, wie häufig eine bestimmte \lstinline|while|-Schleife ausgeführt wird.
    \item ich kann die geeignetste Schleifenart auswählen, um ein bestimmtes Problem zu lösen (\lstinline|for _ in range(zahl)| mit \lstinline|break| oder \lstinline|while|).
    \item Ich kann unterschiedliche Schleifen-Typen mit einem Flussdiagramm modellieren.
    \item Ich kann den \lstinline|return|-Wert einer Funktion direkt und ohne Zwischenspeicherung in einem Ausdruck verwenden, z.B. \lstinline|if rechteck_flaeche(19, 12) < 100: ...|.
\end{todolist}

\section{Lernziele Kapitel 6: Datenstrukturen (Listen, Dictionaries)}


\begin{todolist}
    \item \textbf{Listen: Grundlagen}
    \begin{todolist}
        \item Ich kann Listen in Python erstellen.
        \item Ich kann mit dem Index einzelne Elemente der Liste abrufen und verändern.
        \item Ich kann die Anzahl Elemente einer Liste in Python mit \lstinline|len(liste)| berechnen lassen.
        \item Ich kann die Elemente einer Liste in einer Schleife (\lstinline|for _ in range(zahl)|, \lstinline|while| oder \lstinline|for...in|) durchlaufen.
        \item Ich kann ein Programm schreiben, das in einer Liste die Elemente mit bestimmten Eigenschaften findet (z.B. das Maximum oder alle ungeraden Zahlen).
        \item Ich kann die Elemente einer Liste auf einen Wert reduzieren (z.B. die Summe der Listenelemente berechnen).
        \item Ich kann Listen als Parameter an Funktionen übergeben und innerhalb der Funktion verarbeiten.
    \end{todolist}
    \item \textbf{Algorithmen}
    \begin{todolist}
        \item Ich kann \textit{bubble sort} in Python programmieren.
        \item Ich kann die binäre Suche in Python programmieren.
    \end{todolist}
    \item \textbf{Dynamische Listen}
    \begin{todolist}
        \item Ich kann Listen mit dem Befehl \lstinline|.append(...)| erweitern.
        \item Ich kann Listen mit dem Befehl \lstinline|.pop()| kürzen.
        \item Ich kann Listen mit dem Befehl \lstinline|.insert(position, wert)| an einer bestimmten Position erweitern.
    \end{todolist}
    \item \textbf{Dictionaries}
    \begin{todolist}
        \item Ich kann Dictionaries (Wörterbücher) in Python erstellen und verwenden.
        \item Ich kann die Struktur von Schlüssel-Wert-Paaren in Dictionaries erklären und verstehe den Unterschied zu listenbasierten Datenstrukturen.
        \item Ich kann mit Schlüsseln auf die Werte in einem Dictionary zugreifen.
        \item Ich kann neue Schlüssel-Wert-Paare zu einem Dictionary hinzufügen.
        \item Ich kann bestehende Werte in einem Dictionary ändern.
        \item Ich kann über alle Schlüssel eines Dictionaries mit einer \lstinline|for|-Schleife iterieren.
        \item Ich kann prüfen, ob ein bestimmter Schlüssel in einem Dictionary vorhanden ist (mit \lstinline|if key in dictionary:|).
        \item Ich kann Dictionaries als Parameter an Funktionen übergeben und innerhalb der Funktion verarbeiten.
        \item Ich kann komplexere Datenstrukturen erstellen, indem ich Listen und Dictionaries kombiniere (z.B. Liste von Dictionaries, Dictionary von Listen).
        \item Ich kann alltägliche Anwendungsfälle für Dictionaries identifizieren (z.B. Telefonbuch, Lagerbestand, Preisliste).
        \item Ich kann die Länge (Anzahl der Schlüssel-Wert-Paare) eines Dictionaries mit \lstinline|len(dictionary)| bestimmen.
        \item Ich kann einen Schlüssel-Wert-Paar aus einem Dictionary entfernen.
        \item Ich verstehe, dass Dictionaries ungeordnet sind und die Reihenfolge der Elemente keine Rolle spielt.
    \end{todolist}
\end{todolist}

\section{Lernziele Kapitel 7: Klassen und Objektorientierte Programmierung}
\begin{todolist}
    \item Ich kann erklären, was eine Klasse ist und wie sie sich von einem Objekt unterscheidet.
    \item Ich kann eine einfache Klasse mit Attributen und Methoden definieren.
    \item Ich verstehe den Zweck und die Funktionsweise des Konstruktors (\lstinline|__init__|).
    \item Ich kann erklären, was der Parameter \lstinline|self| in Methoden bedeutet und wie er verwendet wird.
    \item Ich kann Objekte (Instanzen) einer Klasse erstellen und verwenden.
    \item Ich kann auf Attribute und Methoden eines Objekts zugreifen.
    \item Ich verstehe den Unterschied zwischen Klassenattributen und Instanzattributen.
    \item Ich kann das Konzept der Vererbung erklären und einfache Vererbungshierarchien erstellen.
    \item Ich kann Methoden in einer Kindklasse überschreiben und dabei die Methoden der Elternklasse mit \lstinline|super()| aufrufen.
    \item Ich kann reale Probleme mit Hilfe objektorientierter Programmierung modellieren.
    \item Ich kann die Vorteile der objektorientierten Programmierung für komplexe Anwendungen erklären.
\end{todolist}